\chapter{Classification of Reviewed Literature}
\label{app:lit-classification}

As referenced in Chapter 2, this appendix provides a categorical mapping of the key literature reviewed for this thesis. The works are classified based on their primary methodological approach: Model-Based (Electrochemical or Empirical), Data-Driven (RNN, CNN, Transformer), or Hybrid (Physics-Informed).

\begin{table}[h!]
\centering
\renewcommand{\arraystretch}{1.3}
% X column ensures the long category names wrap instead of going off-page
\begin{tabularx}{\textwidth}{@{} l l X @{}}
\toprule
\textbf{Reference} & \textbf{Primary Author} & \textbf{Methodological Category} \\
\midrule

% Model-Based Examples
\cite{naseri_efficient} & Naseri et al. & Model-Based (Electrochemical/DFN) \\
\cite{ma_dual_exp} & Ma et al. & Model-Based (Empirical/Double-Exponential) \\
\cite{verhulst_pso} & Wen/Xian et al. & Model-Based (Empirical/Verhulst) \\

% Data-Driven Examples
\cite{zhang_rnn} & Zhang et al. & Data-Driven (RNN/LSTM) \\
\cite{jafari_cnn_gru} & Jafari et al. & Data-Driven (CNN-GRU Hybrid) \\
\cite{fauzi_sge} & Fauzi et al. & Data-Driven (Transformer/SGEformer) \\

% Hybrid/PINN Examples
\cite{wang_pinn} & Wang et al. & Hybrid (PDE-based PINN) \\
\cite{pde_pinn} & Authors (RG) & Hybrid (PDE-based PINN) \\
\cite{ref1_general} & Review & Comprehensive Review (All Categories) \\

\bottomrule
\end{tabularx}
\caption{Methodological Classification of Cited Works}
\label{tab:lit_class}
\end{table}

\section*{Summary of Classification}
\begin{itemize}
    \item \textbf{Model-Based:} Focuses on interpretability but lacks flexibility for unseen aging patterns.
    \item \textbf{Data-Driven:} Focuses on high accuracy using large datasets but lacks physical consistency.
    \item \textbf{Hybrid:} The target domain of this thesis, aiming to bridge the gap between the two above.
\end{itemize}